\documentclass[11pt,a4paper]{article}
\usepackage[utf8]{inputenc}
\usepackage[T1]{fontenc}
\usepackage{geometry}
\geometry{margin=1in}
\usepackage{hyperref}
\usepackage{enumitem}
\usepackage{booktabs}
\usepackage{longtable}
\usepackage{xcolor}
\usepackage{titlesec}
\usepackage{amsmath,amssymb}

\titleformat{\section}{\Large\bfseries\color{blue!70!black}}{}{0em}{}
\titleformat{\subsection}{\large\bfseries\color{blue!50!black}}{}{0em}{}
\titleformat{\subsubsection}{\normalsize\bfseries\color{blue!30!black}}{}{0em}{}

\title{\textbf{Replication Plan}\\[0.5em]
\large DRAS: Deep Reinforcement Learning for Cluster Scheduling in High Performance Computing\\[0.3em]
\normalsize OSTI ID: 1984484}
\author{Auto-generated Replication Blueprint}
\date{\today}

\begin{document}
\maketitle
\tableofcontents
\newpage

%---------------------------------------------------------------------
\section{Paper Overview}
%---------------------------------------------------------------------
DRAS applies deep reinforcement learning (DRL) to HPC cluster job scheduling. The authors develop CQGym, an OpenAI Gym-compatible environment wrapping CQSim (a cluster queue simulator), and train RL agents (DQL, Policy Gradient, A2C, PPO) to learn scheduling policies that optimize metrics like average bounded slowdown, average wait time, and system utilization. The RL agents are compared against traditional heuristics (FCFS, SJF, WFP, UNICEP) on real HPC job traces.

%---------------------------------------------------------------------
\section{Detailed Workflow}
%---------------------------------------------------------------------

\subsection{Phase 1: Environment Setup}
\begin{enumerate}[label=\textbf{1.\arabic*}]
  \item \textbf{Install CQSim} (cluster queue simulator).
  \item \textbf{Install/build CQGym} (Gym wrapper around CQSim).
  \item \textbf{Install RL libraries}: Stable-Baselines3 or custom implementations of DQL, PG, A2C, PPO.
  \item Verify environment with a test episode.
\end{enumerate}

\subsection{Phase 2: Job Trace Data Preparation}
\begin{enumerate}[label=\textbf{2.\arabic*}]
  \item \textbf{Download HPC job traces.}
    \begin{itemize}
      \item Parallel Workloads Archive: ANL Intrepid, SDSC Blue, RICC, etc.
      \item Convert to CQSim input format (Standard Workload Format).
    \end{itemize}
  \item \textbf{Split traces} into training and test periods.
\end{enumerate}

\subsection{Phase 3: State/Action/Reward Design}
\begin{enumerate}[label=\textbf{3.\arabic*}]
  \item \textbf{State representation.}
    \begin{itemize}
      \item Queue state: number of waiting jobs, resource requests, wait times.
      \item System state: available nodes, running jobs, utilization.
      \item Encode as fixed-size feature vector.
    \end{itemize}
  \item \textbf{Action space.}
    \begin{itemize}
      \item Select which waiting job to schedule next (or select a scheduling heuristic from a portfolio).
    \end{itemize}
  \item \textbf{Reward function.}
    \begin{itemize}
      \item Negative average bounded slowdown, or negative wait time, or utilization-based reward.
    \end{itemize}
\end{enumerate}

\subsection{Phase 4: RL Agent Training}
\begin{enumerate}[label=\textbf{4.\arabic*}]
  \item \textbf{Train DQL, Policy Gradient, A2C, PPO agents} on training traces.
  \item \textbf{Hyperparameter tuning} (learning rate, discount factor, network architecture, replay buffer size).
  \item \textbf{Monitor training curves} (reward per episode, loss).
\end{enumerate}

\subsection{Phase 5: Evaluation and Comparison}
\begin{enumerate}[label=\textbf{5.\arabic*}]
  \item \textbf{Evaluate trained agents} on held-out test traces.
  \item \textbf{Run baseline heuristics} (FCFS, SJF, WFP, UNICEP) on the same traces.
  \item \textbf{Compute metrics}: average bounded slowdown, average wait time, utilization, makespan.
  \item \textbf{Reproduce comparison tables and plots} from the paper.
\end{enumerate}

%---------------------------------------------------------------------
\section{Datasets}
%---------------------------------------------------------------------

\begin{longtable}{p{4.5cm} p{3.5cm} p{6cm}}
\toprule
\textbf{Dataset / Source} & \textbf{Type} & \textbf{Location / Notes} \\
\midrule
\endfirsthead
\toprule
\textbf{Dataset / Source} & \textbf{Type} & \textbf{Location / Notes} \\
\midrule
\endhead
Parallel Workloads Archive & HPC job traces & \url{https://www.cs.huji.ac.il/labs/parallel/workload/} \\
\midrule
ANL Intrepid trace & Job trace & Available from above archive \\
\midrule
SDSC Blue/SP2 trace & Job trace & Available from above archive \\
\midrule
CQSim test traces & Bundled & Included in CQSim repository \\
\bottomrule
\end{longtable}

%---------------------------------------------------------------------
\section{Models, Tools, and Software}
%---------------------------------------------------------------------

\begin{longtable}{p{4.5cm} p{2.5cm} p{6.5cm}}
\toprule
\textbf{Tool / Library} & \textbf{Version} & \textbf{Purpose / URL} \\
\midrule
\endfirsthead
\toprule
\textbf{Tool / Library} & \textbf{Version} & \textbf{Purpose / URL} \\
\midrule
\endhead
CQSim & latest & Cluster queue simulator \newline \url{https://github.com/NSAP/CQSim} (or as referenced) \\
\midrule
CQGym & latest & OpenAI Gym wrapper for CQSim \newline Check paper's supplementary or author's GitHub \\
\midrule
OpenAI Gym / Gymnasium & $\geq 0.26$ & RL environment interface \newline \url{https://gymnasium.farama.org/} \\
\midrule
Stable-Baselines3 & $\geq 1.6$ & RL algorithm implementations (PPO, A2C, DQN) \newline \url{https://stable-baselines3.readthedocs.io/} \\
\midrule
PyTorch & $\geq 1.10$ & Deep learning backend \\
\midrule
NumPy & $\geq 1.21$ & Numerical operations \\
\midrule
Matplotlib / Seaborn & $\geq 3.4$ & Performance comparison plots \\
\midrule
Pandas & $\geq 1.3$ & Job trace data manipulation \\
\bottomrule
\end{longtable}

\subsection{Hardware Requirements}
\begin{itemize}
  \item RL training: GPU recommended but not required ($\sim$hours on CPU).
  \item RAM: $\geq$ 8\,GB.
\end{itemize}

\end{document}
